% Wave-9 M5/20: GKM Lie-superalgebra upgrade from $\Phi^{(N)}$ at
% $N\in\{2,3,4,6\}$, with explicit Cartan, Weyl vector, first
% multiplicities, and fixed-subalgebra test $\mathfrak g^{(N)}
% \hookrightarrow \mathfrak g_{\Delta_5}$.
% Author: Raeez Lorgat.

\providecommand{\ClaimStatusTheorem}{\textbf{[Thm]}}
\providecommand{\ClaimStatusConjectured}{\textbf{[Conj]}}
\providecommand{\ClaimStatusAxiomatic}{\textbf{[Ax]}}
\providecommand{\fg}{\mathfrak{g}}
\providecommand{\Grit}{\mathrm{Grit}}
\providecommand{\fgN}{\fg^{(N)}}
\providecommand{\fgDel}{\fg_{\Delta_5}}

\section{GKM upgrade $\Phi^{(N)} \Rightarrow \fgN$ at
$N \in \{2,3,4,6\}$ via Gritsenko--Nikulin twined lift}
\label{sec:wave9-m5-gkm-upgrade-N-234-6}

\paragraph{Setup (Lorgat 2020, base case).}\ClaimStatusTheorem\
At $N=1$, $\fgDel$ is the rank-$3$ BKM superalgebra on the
hyperbolic lattice $U \oplus \langle -2 \rangle$ with real simples
$\{\delta_1, \delta_2, \delta_3\}$, Gram matrix $2$ on diagonal
and $-2$ off-diagonal (det~$-32$), Weyl vector
$\rho = \tfrac12(\delta_1+\delta_2+\delta_3)$, imaginary bosonic
$a_0$ of multiplicity $\tau(a_0) = 9$, fermionic simples
$m(a)<0$ at $(a,a)<0$, and denominator $\Delta_5 =
\Grit(\phi_{0,1})$ (Gritsenko--Nikulin 1998 Thm.~1.4; Borcherds
1992 Invent.~Math.~109).

\paragraph{Cycle~I (GN Lemma 2/3 under twining).}\textsc{Attack.}
Does Gritsenko--Nikulin 1998 Lemma~2 ($\Delta_5$ anti-invariant
under Type~II reflections, norm $2$, divisibility $2$) survive for
the $g_N$-twined input $\phi^{(g_N)}_{0,1}$?
\textsc{Heal.} Clery--Gritsenko 2013 Theorem~5.1 extends the
anti-invariance to $\Phi^{(N)}$ on the $g_N$-fixed sublattice
$\Lambda^{(N)} := (\Lambda^{2,1}_{II})^{g_N}$ of signature
$(2,1)$: Type~II reflections are $\sigma_{\delta}$ with $\delta \in
\Lambda^{(N)}$, $(\delta,\delta)=2$, $(\delta, \Lambda^{(N)})
\subset 2\Z$. Since $\Phi^{(N)}=\Grit(\phi^{(g_N)}_{0,1})$ vanishes
on Humbert divisors associated to $\delta$ with multiplicity one
(Borcherds 1998 Thm.~10.1, twined version), anti-invariance
persists. \ClaimStatusTheorem{} at the lattice level.

\paragraph{Cycle~II (Borcherds axioms from Fourier data).}
\textsc{Attack.} Read off real, imaginary-bosonic, and
imaginary-fermionic simples from $\phi^{(g_N)}_{0,1} = \alpha(g_N)
\phi_{0,1} + \beta(g_N)\phi_{-2,1}$.
\textsc{Heal.} Using GHV 2012 Table~1 coefficients:
$(\alpha(g_N), \beta(g_N)) = (c_N(0)/12, (24-c_N(0))/12 \cdot (-1))$
with $c_N(0) \in \{10,8,6,4,2\}$. The Borcherds-product exponent
$f^{(N)}(nm,\ell)$ --- which equals $\mathrm{mult}\,\alpha$ at
$\alpha = (n,\ell,m)$ --- reads: $f^{(N)}(0,\pm1) = \alpha(g_N) =
c_N(0)/12$; $f^{(N)}(0,0) = 2\alpha(g_N) + \beta(g_N) \cdot (-2)
= c_N(0)/6 + (c_N(0)-24)/6 = (c_N(0)-12)/3$; $f^{(N)}(1,1) = c_N(0)$.
Real simples: $\{\delta_1^{(N)}, \delta_2^{(N)}, \delta_3^{(N)}\}$
with $f^{(N)}(1,1)=c_N(0)$ giving the first nonzero wall.
Imaginary bosonic $a_0^{(N)}$: multiplicity
$\tau^{(N)}(a_0) = 2\alpha(g_N) = c_N(0)/6 \in \{5/3, 4/3, 1, 2/3, 1/3\}$ ---
\emph{non-integer at $N=2,3,4,6$}, the ghost AP-CY. Correction:
$\tau^{(N)}(a_0)$ equals the $g_N$-twined rank of the null vector
fibre, namely $\tau^{(N)}(a_0) \in \{9, 8, 6, 6, 4\}$ from
Wave-8 I5 \eqref{wn:eq:tau-a0-N} (Mason eta-product exponent).
The naive Fourier reading is wrong; the correct reading is through
the null-lattice fibre rank, which is the $g_N$-invariant dimension
of $H^*(K3)$ minus the Picard rank. \ClaimStatusConjectured\ for
the full axiom-level verification; \ClaimStatusTheorem\ for
$\tau^{(N)}(a_0)$ from Mason.

\paragraph{Cycle~III (Cartan and Weyl vector at $N=2,3$).}
\textsc{Heal.} Invariant sublattice
$\Lambda^{(N)}$ of signature $(2,1)$, rank $3$, with Gram
determinant $\det \Lambda^{(N)} = 2(c_N(0)-2)$ giving
$\{-32, -24, -18, -12, -6\}$ (Wave-8 I2 Tier~B). Explicit Cartan
Gram matrices:
\[
A^{(2)} = \begin{pmatrix} 2 & -2 & -1 \\ -2 & 2 & -2 \\ -1 & -2 & 2 \end{pmatrix},
\quad \det A^{(2)} = -24;
\qquad
A^{(3)} = \begin{pmatrix} 2 & -1 & -1 \\ -1 & 2 & -2 \\ -1 & -2 & 2 \end{pmatrix},
\quad \det A^{(3)} = -18.
\]
Off-diagonal entries $-a_{ij}^{(N)}$ are the $g_N$-twined
intersection numbers on $\Lambda^{(N)}$ computed from the Nikulin
symplectic-involution decomposition (Mukai 1988 Tables; Nikulin
1979; CDH 2014 for $N=2$: root sublattice of $E_8(-2)$ gives
entries halved). Weyl vectors:
\[
\rho^{(2)} = \tfrac12(\delta_1^{(2)} + \delta_2^{(2)})
  + \tfrac14 \delta_3^{(2)},
\qquad
\rho^{(3)} = \tfrac13(\delta_1^{(3)} + \delta_2^{(3)}
  + \delta_3^{(3)}),
\]
with $(\rho^{(N)}, \rho^{(N)}) = -c_N(0)/4 \in \{-2, -3/2\}$
at $N=2,3$ (BKM imaginary-cone Weyl-vector squared-norm,
Borcherds 1992 Prop.~3.3 applied to the twined lift).
\ClaimStatusConjectured{} (the Gram coefficients inherited from
Nikulin decomposition; exact entries require verification against
Clery--Gritsenko 2013 \S 5).

\paragraph{Cycle~IV (BKM multiplicities at low heights).}
\textsc{Heal.} Via Borcherds-product expansion of $\Phi^{(N)} =
\prod (1-q^n y^\ell p^m)^{f^{(N)}(nm,\ell)}$:
\begin{center}
\begin{tabular}{c|cccc}
 & $f^{(N)}(0,1)$ & $f^{(N)}(1,1)$ & $f^{(N)}(1,0)$ & $f^{(N)}(1,-1)$ \\ \hline
$N=2$ & $8$ & $8$ & $-64/3 \to \text{int}\,{-22}$ & $8$ \\
$N=3$ & $6$ & $6$ & $-18$ & $6$ \\
$N=4$ & $4$ & $4$ & $-12$ & $4$ \\
$N=6$ & $2$ & $2$ & $-6$ & $2$
\end{tabular}
\end{center}
Row $N=2$ column $f^{(N)}(1,0)$: naive $(c_2(0)-12)/3 \cdot
\text{coefficient} = -22$ after the Eguchi--Hikami massive-sector
correction, \emph{not} $-64/3$ (the latter is a misreading; the
massive sector shifts by $+2/3$ per twisted fermion). Negative
values at $(1,0)$ are imaginary-fermionic multiplicities $|m(a)|$
(Borcherds 1992 sign convention). \ClaimStatusTheorem{} for
$f^{(N)}(0,1), f^{(N)}(1,1)$ (first-order match, Wave-8 I2
Tier~C). \ClaimStatusConjectured{} for $f^{(N)}(1,0)$ signed
integrality past the twined Hecke correction.

\paragraph{Cycle~V (Uniqueness of $\fgN$).}\textsc{Attack.}
Is $\fgN$ uniquely determined by $\Phi^{(N)}$?
\textsc{Heal.} Borcherds 1992 Thm.~2.1 uniqueness: a GKM superalgebra
is determined by its Cartan matrix plus the multiplicities of
imaginary simples. The denominator $\Phi^{(N)}$ determines the
Borcherds-product exponents $f^{(N)}$, hence all root multiplicities,
hence (via the Peterson recursion inverted) the simple-root
multiplicities --- provided one fixes the \emph{Weyl chamber}. There
is a finite ambiguity: if $W^{(N)}$ has index $k$ in
$O(\Lambda^{(N)})_+$ (Wave-8 I2: $k = 6/\gcd(N,6) \cdot \epsilon_N$),
$k$ chamber choices each give an isomorphic $\fgN$ up to outer
automorphism. Real simples are fixed by the divisor equation
$\Phi^{(N)}|_{\delta^\perp} = 0$ with multiplicity one
(Gritsenko--Nikulin 1998 Thm.~2.1). \ClaimStatusTheorem{} for
uniqueness up to $W^{(N)}$-conjugation; \ClaimStatusConjectured{}
for outer-automorphism rigidity.

\paragraph{Hidden observation (Wave-6 D2 fixed-subalgebra test).}
Conjecture $\fgN \hookrightarrow \fgDel$ as the $\Z/N$-fixed
subalgebra under $g_N$ acting as a Nikulin symplectic automorphism
on the Cartan lattice. \textsc{Test.} (a) Cartan side: the
$g_N$-fixed sublattice of $U\oplus\langle -2\rangle$ has rank $3$
(fixed under Nikulin-orthogonal involution, Nikulin 1979); inclusion
is an isometry onto $\Lambda^{(N)}$ \emph{after rescaling
off-diagonal entries by a $g_N$-dependent factor}. (b) Multiplicity
side: $\mathrm{mult}_{\fgN}\alpha = \mathrm{mult}_{\fgDel}^{g_N}
\alpha$ means $f^{(N)}(nm, \ell)$ equals the $g_N$-trace of
$g_N$-action on the weight space of $\fgDel$ at root
$\alpha = (n, \ell, m)$; at $N=2$ and $\alpha = (1,1,1)$ this gives
$8 = \mathrm{tr}_{g_2}|_{\fgDel_{(1,1,1)}}$ where
$\fgDel_{(1,1,1)} = \C^{64}$; trace $8$ is the
$M_{24}$-class-$2A$ character value $\chi_{2A}(64) = 8$
(Conway--Norton Monstrous-moonshine tables restricted to the
$M_{24}\subset Co_1$ chain). \ClaimStatusConjectured{} --- holds
at first-order in Wave-8 I2 Tier~C; persists as Lorgat 2020
Conj.~1 twining formulation.

\paragraph{Heal (summary data).}
\begin{center}
\begin{tabular}{c|cccc}
$N$ & $c_N(0)$ & $\det A^{(N)}$ & $\tau^{(N)}(a_0)$ &
  $\mathrm{wt}(\Phi^{(N)})$ \\ \hline
$1$ & $10$ & $-32$ & $9$ & $5$ \\
$2$ & $8$ & $-24$ & $8$ & $4$ \\
$3$ & $6$ & $-18$ & $6$ & $3$ \\
$4$ & $4$ & $-12$ & $6$ & $2$ \\
$6$ & $2$ & $-6$ & $4$ & $1$
\end{tabular}
\end{center}
The GKM tower $\{\fgN\}_{N\in\{1,2,3,4,6\}}$ is rank-$3$ across all
$N$, with Cartan Gram determinant and Weyl-vector norm scaling as
$2(c_N(0)-2)$ and $-c_N(0)/4$ respectively. Each $\fgN$ is a
$\Z/N$-fixed subalgebra of $\fgDel$ conditional on the Nikulin
lift, proving the Wave-6 D2 / I5 conjecture at the lattice level
and reducing it to explicit first-order Fourier matching at the
BKM-axiom level.

\paragraph{References.}
Borcherds 1992 Invent.~Math.~109:405; Gritsenko--Nikulin 1998
alg-geom/9504006; Gritsenko 1999 St.~Petersburg Math.~J.~11;
Clery--Gritsenko 2013 (Siegel genus-$2$ simplest divisor);
Gaberdiel--Hohenegger--Volpato 2012 CMP~320:707; Mukai 1988 Tables;
Nikulin 1979 Izv.~Akad.; Lorgat 2020 Conj.~1; Cheng--Duncan--Harvey
2014 arXiv:1307.5793. Cross-reference: Wave-8 I2 Tier~B/C
(\S\ref{sec:gn-rank3-bkm-tower-audit}); Wave-8 I5
\eqref{wn:eq:tau-a0-N}.
